\section{Discussion}\label{sec:discussion}
\subsection{Readiness as scoped evidence}
Readiness is useful when expressed as evidence with scope. A researcher can explain eligibility, a reviewer can identify tested deployment conditions, and an adopter can determine what must be repeated on a different cluster. Explicit incompleteness is actionable: missing distributed evidence calls for another probe or a narrower claim rather than unqualified extrapolation from a local run. The ablation supports this framing. Every conjunct detected faults that a successful smoke run did not reveal (Section~\ref{sec:rq4}), so a smoke run's success should be reported with the obligations it exercised and the plan it described.

The executed analyses also show where the specification needs refinement. Ties at the stability threshold depend on floating-point rounding (Section~\ref{sec:rq1}), and the refinement rule reacts to out-of-memory failures but not to headroom violations that stop short of failure (Section~\ref{sec:rq4}). Both are implementation-level gaps, found by exercising the contract systematically rather than by inspecting individual runs. This supports making such analyses part of the artifact rather than a one-off check.

\subsection{Resource efficiency is comparative}
Resource efficiency is a comparative outcome rather than a property conferred by smoke tests. Validation consumes time and can hold scarce devices. For short reliable workloads, direct execution may be cheaper; for long workloads with detectable early faults, bounded validation may help. A simple break-even condition is $C_{\mathrm{validation}}<pC_{\mathrm{avoidable\ failure}}$, where $p$ is the deployment-specific probability of a failure that the probes actually detect. Neither $p$ nor avoided production cost is estimated by the controller studies. The executed evidence constrains only the left-hand side. Assessment itself is negligible (Section~\ref{sec:rq2}), allocation overhead dominated the pilot's probe computation, and the attempt cap bounds charged validation at $K$ probe allocations. A team's saved allocation also does not establish cluster-wide utilization gains.

The cap is a policy lever with a visible trade-off. A small $K$ bounds cost tightly but prevents refinements that need several halvings from completing (Table~\ref{tab:refinement}); a larger $K$ admits them at proportionally higher validation cost. The choice should follow from the expected cost of a production failure, not from a default.

\subsection{Bounded automation must preserve the research question}
An agent can make a workflow run by changing models, tasks or numerical settings, potentially invalidating the original experiment. The scientific/execution distinction makes that change reviewable, and the refinement study confirms that the implementation never altered $S$. Semantic preservation still needs evidence, and batch invariance is particularly workload-dependent. Our refinement is consequently a constrained proposal, not unrestricted self-repair.

\subsection{Self-application}
An approach for managing HPC experiments itself requires HPC experiments to establish its value. We therefore applied its staged workflow to its own development, progressing from local contract checks to a bounded CUDA integration run and then to one research artifact before considering the twelve-artifact campaign. This makes the proposed discipline concrete: preserve failed or canceled submissions, bind evidence to the configuration actually tested, and require review before expanding the experimental scope. The MapCoder submissions illustrate the point. Three requests were retired before execution, after a configuration defect and unavailable capacity were identified, and each resource-plan change invalidated earlier evidence while the scientific configuration remained fixed. Scheduler admission remains an external constraint; a queued request provides no execution evidence. Self-application can expose implementation and operational defects, but it is not independent validation of effectiveness.

\subsection{Transfer}
Generality depends on transfer conditions. The candidate artifacts vary in workflow, but several depend on cloud APIs or containers, requirements that MCP does not resolve. The likely transferable mechanism is the contract and evidence schema; probes, oracles and feasible resource requests remain artifact- and site-specific. Current evidence supports a reproducible prototype and a path to evaluation, not a claim that twelve diverse studies have already benefited.

\section{Threats to Validity}\label{sec:threats}
\textbf{Construct validity.} Synthetic mutations measure policy behavior on supplied records, not scientific correctness or actual failures prevented. The smoke proxy approximates selected upstream checks. The refinement simulation assumes memory linear in batch size and a fixed per-probe cost, so its convergence counts describe the control loop under that model, not real workloads. Framework-allocated memory excludes some driver and process overhead. The small CUDA network does not approximate long-context inference, agent coordination or training. Retrospective identities in the integration analysis are not pre-execution certificates.

\textbf{Internal validity.} The same developer designed the policy, the operators, the ablation grouping and the oracle, creating favorable-test risk. The ablation removes reason codes after assessment, which is equivalent to removing a conjunct only because every conjunct is evaluated independently. The assessment returns early for nonnumeric telemetry, which the Budget group absorbs. Malformed records may be rejected or raise errors; the prototype has not undergone exhaustive fuzzing. Operator-controlled records can omit costs or assert false coverage. The wrapper is bypassable and has a verification-to-execution race. Inherited completion markers are not content-addressed and can reuse state in ways not fully protected by the wrapper. These limitations prevent security and end-to-end correctness claims.

\textbf{External validity.} One Nibi pilot cannot establish cross-site effects, network behavior or fair-share outcomes. Overhead was timed on one workstation, so absolute latencies will differ elsewhere, although the comparison with probe allocations spans four orders of magnitude. The twelve candidates are purposively selected and concentrated in generation and repair. Current repositories may differ from paper versions, and execution eligibility is pending. API and local-inference workloads have different accounting boundaries. No many-day, multi-node or 64-GPU production experiment is evaluated.

\textbf{Conclusion validity.} Seed variants are not independent studies and receive no inferential tests; deterministic sweeps report exact counts rather than estimates. The planned campaign has a small effective sample and may not establish equivalence. Selecting after outcomes, excluding failed installations or counting successful jobs alone would bias results; the protocol retains such records. Compute savings, long-run reliability and widespread adoption remain hypotheses.

\textbf{Reproducibility.} The execution templates are pinned and attributed. Numbers from the unpublished support manuscript are not reused as new findings. The archive retains controller decisions, ablation, sensitivity, refinement and overhead records, seeds, source hashes, repository-availability records and pilot evidence, and every reported number is regenerated from them. Identity does not capture undeclared dependencies. Release licensing and archival metadata require finalization before public distribution.

\section{Conclusion}\label{sec:conclusion}
\tool{} makes readiness an explicit, bounded decision over identity, coverage and observed resource behavior. The prototype adds assessment and constrained refinement proposals above existing Slurm machinery and exposes them through MCP. The executed studies show that the full contract blocks every constructed fault without false blocks, and that each of its five components is necessary. They also show that decisions follow the declared parameters, that charged validation stays within its cap while proposals leave the scientific specification unchanged, and that assessment costs milliseconds. They also exposed two specification gaps: floating-point ties at the stability threshold and a refinement trigger narrower than the memory conjunct. The GPU pilot exercises limited integration on real hardware. The twelve-artifact protocol identifies the evidence still needed to determine whether these mechanisms save resources while preserving scientific outcomes in real AI software engineering studies.
